% Auto-generated answer document. Original problem statements are preserved verbatim.
\documentclass[11pt]{article}
\usepackage{fontspec}
\usepackage{microtype}
\usepackage{mathtools,amssymb,amsfonts,amsthm,bm}
\usepackage{enumitem}
\usepackage{geometry,fancyhdr,setspace,xcolor}
\usepackage{fancyvrb}
\usepackage{pdfpages}
\usepackage[hidelinks]{hyperref}
\geometry{margin=1in,headheight=14pt}
\setstretch{1.15}
\setlength{\parindent}{0pt}
\setlength{\parskip}{0.55em plus 0.12em minus 0.08em}
\setlength{\abovedisplayskip}{0.85em plus 0.2em minus 0.1em}
\setlength{\belowdisplayskip}{0.85em plus 0.2em minus 0.1em}
\setlist{leftmargin=*,topsep=0.35em,itemsep=0.15em,parsep=0pt}
\pagestyle{fancy}
\fancyhf{}
\fancyhead[L]{\coursename}
\fancyhead[C]{\assignment}
\fancyhead[R]{\studentname}
\fancyfoot[C]{\thepage}
\newcommand{\studentname}{Keshav Ramamurthy}
\newcommand{\coursename}{Math 113}
\newcommand{\assignment}{Homework 1}
\newcommand{\term}{Fall 2026}
\newcommand{\R}{\mathbb{R}}
\newcommand{\C}{\mathbb{C}}
\newcommand{\Q}{\mathbb{Q}}
\newcommand{\Z}{\mathbb{Z}}
\newcommand{\N}{\mathbb{N}}
\newcommand{\F}{\mathbb{F}}
\newcommand{\K}{\mathbb{K}}
\newcommand{\E}{\mathbb{E}}
\newcommand{\Prob}{\mathbb{P}}
\newcommand{\abs}[1]{\left\lvert#1\right\rvert}
\newcommand{\norm}[1]{\left\lVert#1\right\rVert}
\newcommand{\ip}[2]{\left\langle#1,#2\right\rangle}
\newcommand{\set}[1]{\left\{#1\right\}}
\newcommand{\given}{\,\middle\vert\,}
\newcommand{\dd}{\,\mathrm{d}}
\DeclareMathOperator{\rank}{rank}
\DeclareMathOperator{\nullity}{nullity}
\DeclareMathOperator{\Span}{span}
\DeclareMathOperator{\Ker}{ker}
\DeclareMathOperator{\im}{im}
\DeclareMathOperator{\Hom}{Hom}
\DeclareMathOperator{\Aut}{Aut}
\DeclareMathOperator{\Var}{Var}
\DeclareMathOperator{\Cov}{Cov}
\newtheorem{problem}{Problem}[section]
\newenvironment{solution}{%
  \par\smallskip\noindent\textbf{Solution notes}\par
  \vspace{1.15in}\noindent\hrulefill\par\vspace{0.28in}%
}{\par}
\newenvironment{originalproblem}{%
  \VerbatimEnvironment
  \par\medskip\noindent\begin{minipage}{0.98\linewidth}
  \begin{Verbatim}[fontsize=\small]%
}{%
  \end{Verbatim}\end{minipage}\par
}
\renewcommand{\coursename}{Math 113}
\renewcommand{\assignment}{Homework 1 Answers}
\begin{document}
\begin{center}\Large\textbf{Math 113 -- Homework 1 Answers}\end{center}
\includepdf[pages=1,pagecommand={\thispagestyle{empty}}]{hw01.pdf}
\clearpage
\section{Expanded problems for credit}
\newpage \subsection*{Problem 1}
\begin{originalproblem}
\end{originalproblem}
\begin{solution}
    We consider what this set is. \(\{n \in \mathbb{Z} \mid n^2 < 0\}\) describes
all integers where $n^2 < 0.$\par Notice by the trivial inequality $x^2 \ge 0 $ for $x \in \mathbb{R}$.
Thus, no integers satisfy this condition and it describes the empty set. 
% this, this relation creates a partition of $\mathbb{R}$, 
% Write your proof or calculation here.
\end{solution}
\newpage \subsection*{Problem 2}
\begin{originalproblem}
\end{originalproblem}
\begin{solution}
\begin{enumerate}
    \item The empty set has only one subset, $\emptyset$. Thus, it's power set is
        $\{\emptyset\}$ with cardinality 1
    \item The power set of $\{a\}$ is $\{\emptyset, \{a\}\}$ so its cardinality is
        2.
    \item The power set of $\{a, b\}$ is $\{\emptyset, \{a\}, \{b\}, \{a, b\}\}$ with cardinality
        4
    \item  The power set of $\{a, b, c\} $ is $\{\emptyset, \{a\}, \{b\}, \{c\},
        \{a, b\}, \{a, c\}, \{b, c\}, \{a, b, c\}\}$ with cardinality 8
\end{enumerate}
% Write your proof or calculation here.
\end{solution}
\newpage \subsection*{Problem 3}
\begin{originalproblem}
\end{originalproblem}
\begin{solution}
For the pair $(0, 3) \in \mathbb{R} \times \mathbb{R}$ for $\ge$ to be an equivalence
relation, the symmetric property would have to hold, but notice $0 \not\ge 3, 3 \ge 0$, so it
can't be an equivalence relation.
% Write your proof or calculation here.
\end{solution}
\newpage \subsection*{Problem 4}
\begin{originalproblem}
\end{originalproblem}
\begin{solution}
    We have $(a * b) * c = (b * c) = a.$ \par For the other, we have $a * (b * c) = a * a =
    a$. \par You can not say on the basis of this one calculation, it would have to hold over
    all possible values associativity holds.
% Write your proof or calculation here.
\end{solution}
\newpage \subsection*{Problem 5}
\begin{originalproblem}
\end{originalproblem}

\begin{solution}
We have $a * b = ab + 1$ and $b * a = ba + 1$, so $a * b = b * a$, and commutativity
holds.
\par We have $$a * (b * c) = a * (bc + 1) = a(bc + 1) + 1 = abc + a + 1$$ and we have $$(a
* b) * c = (a * b + 1) * c = (a * b + 1)c + 1 = abc + c + 1$$which are not equal when
$a \ne c$, so it is not associative.
% Write your proof or calculation here.
\end{solution}
\newpage \subsection*{Problem 6}
\begin{originalproblem}
\end{originalproblem}
\begin{solution}
% Write your proof or calculation here.
We first find $\phi(a + b) = -(a + b) = -a - b = \phi(a) + \phi(b)$ which establishes
$\phi$ as a homomorphism
\par To prove that the function $\phi(a) = -a$ is injective, we see that $$\phi(a) =
\phi(b) \iff -a = -b \iff a = b$$
\par To prove surjectivity, we see that $$\forall x \in \mathbb{Z},  \phi(-x)\in \mathbb{Z} \text{ and } \phi(-x) = x$$ and if
$x \in \mathbb{Z} \iff -x \in \mathbb{Z}$
Thus, it satisfies the homomorphism condition, injectivity, surjectivity making it an
isomorphism
\end{solution}
\newpage \subsection*{Problem 7}
\begin{originalproblem}
\end{originalproblem}
\begin{solution}
Suppose we have a binary structure $(S, *)$ and there is an isomorphism $\phi : S
\rightarrow S' $ mapping $(S, *)$ to $(S', {*}^{'})$. %\par For sake of contradiction, let's assume that the binary operator
% ${*}^{'}$ does not have associativity.
\par Subject to the conditions of homomorphism, we have that $\phi(a * b) = \phi(a) *^{'}
\phi(b)$, and along with this, we're given the condition of associativity and
homomorphism for $(S, *)$
which yields $$(a * b) * c = a * (b * c) \rightarrow \phi((a * b) * c) = \phi(a * (b *
c))$$
The LHS is 
$$\phi((a * b) * c) = ((\phi(a) * \phi(b)) *^{'} \phi(c) $$
and the RHS is 
$$\phi(a * (b * c)) = \phi(a) *^{'} (\phi(b) *  \phi(c))$$
 We now apply the homomorphism again on two variables, namely $\phi(a * b) = \phi(a)
*^{'} \phi(b)$
\par Substituting yields 
$$(\phi(a) *^{'}\phi(b))*^{'} \phi(c)= \phi(a) *^{'} (\phi(b) *^{'} \phi(c))$$
which are equal, because the original LHS and RHS were identical
\par To conclude, the last property of isomorphism is surjectivity, and so notice for arbitrary $$a', b', c'
\in S^{'}, \exists a, b, c \in S : \phi(a) = a', \phi(b) = b', \phi(c) = c'$$
and so we can substitute the equation above using this mapping.
$$(x' *^{'} y') *^{'} z' = x' *^{'} (y^{'} *^{'} z^{'}) \text{ proving associativity} \qed$$
% Write your proof or calculation here.
\end{solution}
\section{Practice problems}
\begin{originalproblem}
Practice problems (not for submission):
Fraleigh Exercises 0.5, 0.9, 0.11, 0.17, 0.25, 0.28, 0.29, 0.31, 0.36;
2.3, 2.17, 2.26, 2.28, 2.37; and 3.3, 3.4, 3.24, 3.25, 3.31, 3.33.
\end{originalproblem}
\end{document}
